\clearpage
\fancyhead[L]{\small PASQUICH. STERNWARTE IN OFEN.}
\hypertarget{native-502-lower}{}\label{native-502-lower}
\begin{center}
\small Göttingische gelehrte Anzeigen. No. 85. pp. 841..843. 28 May 1808.\\[3pt]
\rule{.72\linewidth}{.3pt}
\end{center}

Under the title \emph{Account of my proposals for advancing astronomy at the Royal University Observatory at Ofen}, the director of the observatory, Mr \textsc{Joh. Pasquich}, has had a short pamphlet printed. It gives very welcome grounds for hoping that astronomy will flourish in Hungary.

{\rightskip=0pt plus 3em The observatory, which has stood at Ofen since 1780, is unsuited to the present state of astronomy both in its construction and in its equipment. This was all the more regrettable because the university statutes require the observatory to maintain two assistants and a caretaker in addition to the director. The Royal Hungarian Lieutenancy, headed by the enlightened Archduke Palatine, recognized how important the establishment could become for science and for the country if men who were active and wholly devoted to their science were given better equipment there. Professor \textsc{Pasquich} therefore received an order from the highest authority, immediately upon his appointment as astronomer at the observatory, to submit proposals for procuring new instruments. He did so on the principle that the real aim was the genuine advancement of science, to which all cultured nations of Europe consider themselves bound to contribute. The Archduke Palatine and the Royal Lieutenancy, as well as the Emperor himself, approved his proposals immediately and without reservation, and ordered that the instruments be commissioned without delay. He therefore approached artillery captain \textsc{Reichenbach} in Munich, who, as is well known, set up a workshop there some years ago for making mathematical instruments and, judging by several trials, seems able to rival the foremost English instrument makers. The order was for a three-foot repeating circle with a thirty-inch azimuth circle and a silver graduated limb, a six-foot meridian telescope, an astronomical seconds-pendulum clock, an astronomical travelling clock with a half-second pendulum, an eighteen-inch
\SourceLeaf{503}
astronomical circle, a twelve-inch circle for terrestrial measurements, an eight-foot telescope, and an equatorial. All this, including some repairs to other instruments, was agreed at the moderate price of 7210 imperial gulden or 8652 gulden in Reich currency. In addition, an astronomical seconds-pendulum clock was ordered from mining councillor \textsc{Seyffert} in Dresden for 360 thalers. The war interrupted the making of those instruments, but the work is now complete except for the eight-foot telescope and the equatorial. Some of the instruments have already arrived at Ofen, and the rest are expected any day.\par}

All that remains to be wished for now is that this fine equipment may soon find a home worthy of it, where it can be used for the benefit of science and to the honour of the Hungarian nation and its distinguished patrons. Mr \textsc{Pasquich} has already proposed a suitable site for the new observatory. Only a few special circumstances leave it uncertain whether that site will be chosen. In any event, one can confidently expect that Hungary will soon be graced with a temple of Urania, an achievement on which one will be able to congratulate the Hungarian nation and science itself.
